\rhead{CS301: Computer Networks}
\phantomsection 
\section*{Computer Networks (CS301)}
\label{course:CS301}

\noindent \small{\hyperref[main:scheme]{$\leftarrow$ Back to Scheme}} \vspace{0.5cm}

\noindent \textbf{Credits:} 3 (3-0-0)

\subsection*{Course Content}
\noindent\textbf{Unit 1} \\
\textit{Network Fundamentals and Reference Models:} Network types (LAN, MAN, WAN, PAN) and topologies, Client-server vs. peer-to-peer architectures, OSI vs. TCP/IP models with layer responsibilities, Encapsulation/decapsulation process, Protocol data units (PDU, frame, packet, segment), Network devices (hubs, switches, routers, gateways), Performance metrics (bandwidth, throughput, latency, jitter, goodput), Interview questions on layered architectures and protocol interactions.

\vspace{0.3cm}

\noindent\textbf{Unit 2} \\
\textit{Physical and Data Link Layers:} Transmission media (twisted pair, coaxial, fiber optic, wireless), Signal encoding (NRZ, Manchester, 4B/5B), Error detection (parity, checksum, CRC polynomials), Error correction (Hamming codes), Framing methods (byte/packet/character stuffing, bit-oriented), Flow control (stop-and-wait, sliding window), ARQ protocols (Go-Back-N, Selective Repeat), MAC protocols (ALOHA, CSMA/CD, CSMA/CA).

\vspace{0.3cm}

\noindent\textbf{Unit 3} \\
\textit{Network Layer - Addressing and Routing:} IPv4 addressing (classes, CIDR, subnetting, VLSM, supernetting), IPv6 addressing and transition mechanisms, ARP/RARP protocol operation, ICMP error reporting and diagnostics, Routing algorithms (distance vector, link state, hierarchical), Routing protocols (RIP, OSPF, BGP), Static vs. dynamic routing, Route summarization and longest prefix matching.

\vspace{0.3cm}

\noindent\textbf{Unit 4} \\
\textit{Network Layer - Congestion and Internetworking:} Congestion control mechanisms (leaky bucket, token bucket), Quality of Service (QoS) concepts (DiffServ, IntServ), NAT/PAT operation and port address translation, Fragmentation and reassembly, MPLS basics and label switching, Internetworking challenges (heterogeneous networks, fragmentation, MTU issues), Interview problems on subnet calculation, CIDR aggregation, and traceroute analysis.

\vspace{0.3cm}

\noindent\textbf{Unit 5} \\
\textit{Transport and Application Layers:} UDP protocol characteristics and use cases, TCP protocol (3-way handshake, state diagram, connection release), TCP flow control (sliding window, silly window syndrome), TCP congestion control (slow start, congestion avoidance, fast retransmit/recovery), Reliable data transfer principles, Application layer protocols (DNS resolution, HTTP/HTTPS, SMTP, FTP, DHCP), Socket programming fundamentals and port number assignments.

\newpage
\rhead{Curriculum Schema}
